\subsection{Classical coding and the precise forced model}
\label{sec:related-work}

The literature supplies several of the mechanisms used here. Write
$T(2k)=k$ and $T(2k+1)=3k+2$ for the shortcut map. Lagarias's
survey \cite{lagarias-1985} places parity coding, stopping times, inverse
trees, and generalized maps in their historical setting. For the ordinary
shortcut map, Terras \cite{terras-1976} proves the affine dependence of an
iterate on its initial value and the correspondence between finite parity
words and residue classes. Bernstein's explicit inverse coding
\cite{bernstein-1994} and the conjugacy developed by Bernstein and
Lagarias \cite{bernstein-lagarias-1996} describe the associated $2$-adic
dynamics. Proposition~\ref{prop:affine-parity} uses this classical
mechanism with a time-dependent intercept. Its contribution is the
explicit forced formula and its role in the subsequent certificates.

The companion generator in our height coordinate also has an exact
precedent. With $P(x)=x+1$, the map $U=PTP^{-1}$ is the map $E$ of
Monks et al.\ \cite[Section~3]{monks-et-al-2013}. Barina uses the same
translated map in an ordinary-trajectory batching algorithm
\cite[Section~3.1]{barina-2025}. The height proposition here establishes
which externally prescribed impulse words admit a binary height,
identifies their safe domain, and produces the resulting periodic
schedule of $T$ and $U$. The translated generator itself is inherited.

There are two distinct uses of periodicity in this subject. A map may
have coefficients periodic in the residue of its \emph{input}; our
impulses are periodic in \emph{time}, so the state includes the phase.
For a word of period $p$, a full return at a fixed phase is
residue-affine:
\[
 G(n)=\frac{3^{o_r}n+b_r}{2^p},\qquad n\equiv r\pmod{2^p}.
\]
This fits the independent-offset generalized class stated by Kionke
\cite[Section~1]{kionke-2017}. It should not automatically be identified
with every earlier definition of a generalized Syracuse map. In
particular, the original definition of Matthews and Watts
\cite[p.~167, equation~(1.3)]{matthews-watts-1984} uses a common system
of residue representatives, an additional restriction that need not
hold for a forced return. We use that source only for this checked
definition, rather than invoking a theorem from an uninspected portion
of its text.

The locally scaling shift theorem of Kingsbery et al.
\cite[Theorem~4.8 of the preprint]{kingsbery-et-al-2011} supplies a
further $2$-adic framework for these returns. Such coding concerns the
$2$-adic space, where a valid infinite itinerary need not represent a
positive integer. Akin's analysis of exceptional rational points
\cite{akin-2004} makes this distinction particularly relevant. Neither
a symbolic conjugacy nor the existence of every finite parity prefix
provides the real-size descent needed for integer capture.

\subsection{Stopping-time results and universal capture}
\label{sec:density-literature}

The natural-density first-descent theorem in Terras
\cite{terras-1976} and the logarithmic-density almost-bounded-orbit
theorem of Tao \cite[Theorem~1.3]{tao} have different conclusions from
capture of every seed. Tao allows any prescribed bound tending to
infinity and proves that the orbit minimum is below that bound outside
a set of logarithmic density zero. The possible exceptional set remains
material here: our conditional counting argument requires an upper
bound $o(X^\delta)$ for the exceptional set of the \emph{same forced
map}. Even a natural-density-one capture result would give only $o(X)$
exceptions, which does not supply that stronger comparison.

The intermediate thresholds also matter. Korec
\cite[Theorem~1]{korec-1994} obtains natural-density-one descent below
$n^c$ for each fixed $c>\log_4 3$. Inselmann's preprint
\cite[Corollary~1.4, arXiv v3]{inselmann-2024} gives, for every fixed
$\varepsilon>0$, the bound
$T^{\lfloor 2\log n/\log(4/3)\rfloor}(n)\leq n^\varepsilon$
on a set of natural density one. These ordinary-map statements distinguish
the size of the target, the density notion, and the iteration horizon.
Their transfer to our forced map would require an additional argument.
Shaik's recent preprint \cite[Theorem~1.1 and Section~8]{shaik-2026}
reports natural-density descent to a polylogarithmic target with an
arbitrary divergent multiplier, together with time and orbit-height
bounds. Its author reports a Lean proof of the main theorem and explicitly
distinguishes some registered corollary normalizations from the printed
ones. We have not independently rebuilt or axiom-audited that artifact;
these reported ordinary-map results are comparisons, not proof inputs.

For independent residue slopes, Venturini's published abstract
\cite{venturini-1997} states natural-density-one fixed-factor descent
in absolute value on each signed half-line, under positive coprime
branch numerators and geometric mean slope below one. Our period-$p$
return satisfies these stated algebraic conditions: the parity bijection
gives the slope product $(3/4)^{p2^{p-1}}<1$. The abstract allows a fixed
bounded alternative, which is absorbed by any fixed-factor bound for
sufficiently large inputs. The full proof was not recovered in this
review. This is a direct model comparison, but its conclusion supplies
neither capture nor descent below every prescribed divergent threshold.

Gon\c{c}alves, Greenfeld and Madrid
\cite[Definition~1.2 and Theorems~1.3 and~1.8 of arXiv v2]{goncalves-greenfeld-madrid-2025}
extend almost-bounded results to a
specified family of generalized Collatz maps. Their principal family
uses a fixed multiplier $q$ on the nonzero residue classes modulo a
fixed $p$, with positivity, integrality, coprimality, and the inequality
$q<p^{p/(p-1)}$; their natural-density power conclusion has additional
positive-offset hypotheses. Our full return generally has the different
multipliers $3^{o_r}$. Checking an average slope is therefore insufficient
to invoke their theorem. Applying it to this forced family would require
a proved reduction to their exact model.

Siegel and O'Dwyer's recent preprint
\cite[Theorem~5 and Definition~19, arXiv v1]{siegel-odwyer-2026}
reports an almost-boundedness theorem on number rings with varying
slopes. Its intended scalar conditions are compatible with our seven-step
return: $\rho_H=(3/4)^{448}$ and $127^2\rho_H<1$. This makes it a close
comparison, but we do not invoke its claimed theorem. The literal
global branch-integrality requirement in Definition~16 conflicts with
the coset-only condition in Definition~19, and the passage-probability
argument at the end of Section~7.2 requires further justification.
These are unresolved concerns about the inspected version, not a
counterexample to its intended conclusion. None of our capture
hypotheses is discharged by that preprint.

Kionke also states a one-dimensional eventual-periodicity conjecture for
coprime residue-affine maps whose product of absolute slope numerators
is smaller than the appropriate power of their denominator
\cite[Section~1]{kionke-2017}. That statement is a conjecture, not a
capture theorem available to close our residual families. It quantifies
over all signed integers. Even a future proof for every positive seed
of our fixed word would require a separate negative-domain analysis
before it could settle an instance in that form. The higher-dimensional
counterexamples in the same paper do not settle our one-dimensional
problem either.

\subsection{Inverse growth and arithmetic progressions}

Applegate and Lagarias \cite{applegate-lagarias-1995-i} develop finite
inverse-tree enumeration, including odd-step and congruence information,
as a way to obtain infinite lower counting bounds. Their second paper
\cite{applegate-lagarias-1995-ii} develops the difference-inequality
approach further. Krasikov and Lagarias
\cite[Theorems~2.2 and~6.1]{krasikov-lagarias} provide a soundness
framework and an ordinary-map lower exponent $0.84$ for positive targets
coprime to three. Our exponent between $0.535$ and $0.536$ concerns a
different, forced map and is not a numerical improvement to that result.
The model-specific work is the complete eligible branch census,
the forced affine offsets, distinctness across depths, the stopping
frontier, and the transfer to all safe targets and progressions.

The progression comparison is especially close. Monks
\cite{monks-2006} proves that every arithmetic progression is
sufficient for ordinary Collatz, in the sense of eventual merger with
an orbit from the progression. Monks et al.\ \cite{monks-et-al-2013}
develop admissible backtracing and stronger sufficiency statements.
These are direct antecedents of the affine residue steering used here.
Sufficiency does not assert that every forward orbit visits a given
progression. Their bounded odd-inverse count for targets coprime to
three also differs from the literal selector length needed here;
Section~\ref{sec:inverse} proves the size, multiplicity, and attachment
estimates required by our counting transfer.

A possible quantitative improvement is to retain more phase and residue
information in a system of forced inverse inequalities. The classical
inequality method suggests this experiment, but its soundness proof
would need the actual shifted arguments, eligibility guards, and
coverage of our return. Importing an ordinary-map recurrence or optimizing
unverified numerical weights would not establish a new exponent.

\subsection{Common control and synchronization}

Single-state steering, a common word for a finite ensemble, and periodic
capture from every phase are separate levels of control. Vasilyev's
recent preprint \cite[Section~4]{vasilyev-2026} studies the same available
upward and downward integer edges and gives closely related
single-state steering. A state-dependent choice of an edge does not
produce one temporal word that moves several states simultaneously.
We use a different notation for the translated companion $U$, so the
symbol $U$ in that preprint must not be identified with ours.

The pair-merging induction is classical in finite automata:
\v{C}ern\'y's criterion \cite[Theorem~2]{cerny-1964} reduces
synchronization to merging every pair. The arithmetic merger and
protected reference make that induction useful in our infinite state
space. The logarithmic dependence on the interval endpoint, with a word
chosen for that interval, and success of every cyclic rotation are
additional conclusions proved here, not consequences of the
finite-state criterion alone.

The probabilistic proof has the structure of drift and minorization
arguments used for geometric convergence \cite{rosenthal-1995}.
Shared random maps also occur in coupling from the past
\cite{propp-wilson-1996}, although our objective is word existence,
not exact sampling from a stationary distribution. Yano and Yasutomi
\cite{yano-yasutomi-2011} construct synchronizing random-map
realizations for finite ergodic chains. Here the maps and their common
coupling are prescribed, and the ensemble is not a finite reference
chain. We therefore prove the required ensemble estimates, reference
closure, and all-rotation union bound rather than infer them from
finite-chain mixing.

\subsection{Continuous geometry and inherited attraction}

The cosine continuation is Chamberland's function
\cite{chamberland-1996,chamberland-survey}; its interpolation is an
explicit modeling choice. Dumont and Reiter
\cite{dumont-reiter-2001,dumont-reiter-2003} already study real and
complex visualizations and alternative continuations. Agreement on
integer inputs leaves substantial freedom between them. In particular,
results about a chosen holomorphic extension, such as those of
Letherman, Schleicher and Wood \cite{letherman-schleicher-wood-1999},
cannot be transferred solely from their common integer values.

The general integer-cycle attraction proposition has a stronger prior
relationship than analogy. After the positive-height coordinate change,
the full-period return satisfies the hypotheses of Sindee, Sa
Ngiamsunthorn and Sriwongsa \cite[Theorem~3.3]{sindee-2024}.
Their extension and the positive-coordinate conjugate of our composed
cosine return have the same values and first derivatives at integers.
Section~\ref{sec:geometry} gives the
reduction and retains a direct multiplier proof. The general attraction
mechanism is therefore attributed as an application of existing work.
The particular trapping disks, noninteger cycles, real separator, and
phase-dependent complex trajectory still require the separate
certificates supplied here.

Lygeros and Rozier \cite{lygeros-rozier-2014} give shrinking real
intervals around integers for the unforced cosine map. A forced analogue
would be a useful next geometric theorem: specify interval widths at
each phase, prove inclusion under the actual translated update, and
handle the finitely many small centers. The original unforced
inequalities do not provide that inclusion after adding impulses.

\subsection{Finite certificates, termination, and formal proof}

The unresolved seven-step families can be viewed as a termination
problem on exact arithmetic domains. The size-change principle of
Lee, Jones and Ben-Amram \cite{lee-jones-benamram-2001}, and the
corrected development of monotonicity constraints by Ben-Amram
\cite[Section~3]{benamram-2010-monotonicity}, suggest finite graphs
whose edges certify comparisons of well-founded counters. To use this
method, each edge must hold on its entire guarded domain, every possible
execution must be covered, and the composed graph must force an infinite
descending thread along any putative infinite execution. Finding a
cycle in a residue graph, or checking only sampled parameters, does not
establish these conditions.

Transition invariants \cite{podelski-rybalchenko-2004} provide another
relevant design principle: a finite union of well-founded relations can
cover the transitive closure of the transition relation. Coverage of
execution segments is the substantive obligation. Path-complete graph
Lyapunov functions \cite{ahmadi-et-al-2014} similarly clarify why all
allowed products must be represented. Their matrix-system hypotheses
do not automatically cover guarded affine integer returns. An affine
lift must also account for offsets and the extra constant coordinate;
unconstrained matrix contraction is not supplied by a list of favorable
branch slopes.

Yolcu, Aaronson and Heule \cite[Section~3.2]{yolcu-aaronson-heule-2023}
give a mixed-base string-rewriting encoding whose termination is
equivalent to ordinary Collatz. Their negative results concern specified
classes of interpretations, not all termination methods. A forced
encoding would need its own phase representation, forward and backward
simulation, and proof that administrative rewriting cannot continue
forever without physical progress. Carelli's connection between
one-variable linear-constraint loops and weak generalized Collatz
reachability \cite{carelli-2026} is also conditional on an explicit
conjecture and an exact model. It is not an automatic decision procedure
for our guarded, multi-branch return.

Nashida's recent companion preprints
\cite{nashida-2026,nashida-2026-part-ii} address staged valuation
certificates and limitations of finite abstractions and matrix
interpretations. Their revision-r3 repositories report Lean coverage
of certificate soundness and five examples in Part~I, and seven
residue-avoidance weakenings in Part~II. The general dichotomy,
nonexistence results, matrix barriers, and maximality are outside the
reported formalization. We inspected those scope statements without
independently rebuilding or axiom-auditing the artifacts. A use for our
forced system would need its own faithful encoding and certificate
verification; the reported barriers concern their specified methods.

For certificate formalization, Gonthier's computational-reflection
architecture \cite{gonthier-2005} illustrates the essential separation:
prove a checker's semantic correctness, then use its accepted finite
data to obtain the mathematical statement. Solovyev and Hales
\cite{solovyev-hales-2013} provide an analytic example using Taylor
interval certificates. Our next Lean milestones are the corresponding
coverage and soundness theorems, not merely imports of recorded maxima
or successful external runs. The Lean infrastructure and its axiom and
validation guidance \cite{demoura-ullrich-2021,lean-axioms-424,
lean-validating-proofs} help describe that trust boundary; the exact
coverage of this repository is listed in Section~\ref{sec:formal-scope}.

Finally, the residue-affine literature gives useful limits on overly
simple descriptions. Kohl proves wildness of iterates for surjective
noninjective residue-class-wise affine maps \cite{kohl-2007}, and a
separate obstruction for a restricted class of monotonizing conjugacies
\cite{kohl-2008-conjugates}. These statements do not rule out
variable-time returns, relational mergers, residue-dependent ranks, or
finite certificates of entry. They support specifying the exact scope
of an obstruction before using it to discard a proof strategy.

\subsection{Contribution and review boundaries}
\label{sec:literature-boundaries}

The release combines a model-specific admissibility classification,
exact support bounds, an all-phase synchronization theorem, forced
inverse growth and progression transfer, explicit geometric
certificates, and conditional seven-step reductions. Their classical
coding, control, counting, interpolation, and verification ingredients
are attributed above and beside their uses. No general novelty claim
is made for those ingredients or for the inherited integer-attraction
principle. Priority for the precise forced statements remains a separate
review question.

The bibliography is selective, rather than a claim to enumerate every
Collatz paper. The accompanying staging audit accounts for all 88 entries
in the pre-existing source register, records additional sources, and
distinguishes inspected primary text from metadata or abstract-only
checks. In particular, the Matthews--Watts comparison above uses the
recovered definition; the Letherman--Schleicher--Wood comparison uses
the publisher abstract and the explicit distinction in Dumont--Reiter;
the Lee--Jones--Ben-Amram and Podelski--Rybalchenko references supply
historical and methodological context whose exact formal use would
require further source review. Full-text gaps and version discrepancies
remain visible in that audit. None is silently treated as a theorem
that closes the remaining capture hypotheses.
